% GENERATED FILE. Edit canonical lessons, then run npm run generate:books.
% canonical-source-hash: fnv1a64:152947efa39a51bc
% canonical-lessons: MR-C366-mulya, MR-C366-naitikta, MR-C366-sacoti, MR-C366-tyag, MR-C366-hit

\chapter{Values}
\label{ch:mr-values}
\hlchaptermodality{366}

\begin{chapteropening}
\textbf{By the end of this chapter:} \emph{I can talk about values, integrity and sacrifice.}

Complete the last lesson of chapter 366: I can talk about values, integrity and sacrifice.
\end{chapteropening}

\section[\texorpdfstring{\mr{मूल्य} (mūlya)}{मूल्य (mūlya)}]{\mr{मूल्य} (mūlya) — a value}

\label{lesson:MR-C366-mulya}



\begin{quote}
Before the new one: say the Marathi for to resolve, then the Marathi for to fight for.
\end{quote}

\subsection*{You'll want to know: \mr{मूल्य}}
\textbf{\mr{मूल्य}} — \emph{mūlya} — \textquotedblleft{}a value\textquotedblright{}.

\begin{quote}
\mr{आईवडिलांनी} \mr{आम्हाला} \mr{चांगली} \mr{मूल्यं} \mr{शिकवली}.
\end{quote}

\begin{quote}
\emph{āīvḍilānnī āmhālā cāṅglī mūlyaṁ śikavlī.}
\end{quote}

\begin{quote}
Our parents taught us good values.
\end{quote}

Neuter. Plural: \mr{मूल्यं}, \emph{mūlyaṁ}.

Also price, in formal use: \mr{वस्तूचं} \mr{मूल्य}.

\begin{cousinweb}
\textbf{\mr{मूल्य}} (\emph{mūlya}) is a tatsama, taken from Sanskrit \emph{mūlya}, \textquotedblleft{}original value, price\textquotedblright{} (CDIAL 10257).
\end{cousinweb}

\begin{grammarlens}[title={What you've built}]
One more word for giving reasons and opinions.
\end{grammarlens}

\subsection*{Your turn}
\begin{itemize}
\raggedright
  \item \emph{Say it:} \emph{mūlya}
  \item \emph{Say it:} \emph{mūlya}, once more
  \item \emph{Say it:} \emph{āīvḍilānnī āmhālā cāṅglī mūlyaṁ śikavlī.} — Our parents taught us good values.
  \item \emph{Recall:} say the Marathi for to resolve, then the Marathi for to fight for, then say \emph{mūlya} again
\end{itemize}

\begin{tcolorbox}[breakable,colback=teal!4,colframe=teal!35!black,title={Before you move on}]
What does \mr{मूल्य} mean? (\textquotedblleft{}A value\textquotedblright{}.) Say it once more.
\end{tcolorbox}
\section[\texorpdfstring{\mr{नैतिकता} (naitiktā)}{नैतिकता (naitiktā)}]{\mr{नैतिकता} (naitiktā) — morality; ethics}

\label{lesson:MR-C366-naitikta}



\begin{quote}
Before the new one: say the Marathi for to fight for, then the Marathi for a value.
\end{quote}

\subsection*{You'll want to know: \mr{नैतिकता}}
\textbf{\mr{नैतिकता}} — \emph{naitiktā} — \textquotedblleft{}morality; ethics\textquotedblright{}.

\begin{quote}
\mr{व्यापारातही} \mr{नैतिकता} \mr{पाळली} \mr{पाहिजे}.
\end{quote}

\begin{quote}
\emph{vyāpārāthī naitiktā pāḷlī pāhije.}
\end{quote}

\begin{quote}
Ethics should be observed in trade too.
\end{quote}

Feminine.

\mr{नैतिक}, moral, + -\mr{ता}.

\begin{cousinweb}
\textbf{\mr{नैतिकता}} (\emph{naitiktā}) is a learned word built on Sanskrit \emph{nīti}, \textquotedblleft{}conduct, guidance\textquotedblright{} (CDIAL 7546), through the adjective \mr{नैतिक}, \textquotedblleft{}moral\textquotedblright{}, with \emph{-tā}, \textquotedblleft{}-ness\textquotedblright{}.
\end{cousinweb}

\begin{grammarlens}[title={What you've built}]
One more word for giving reasons and opinions.
\end{grammarlens}

\subsection*{Your turn}
\begin{itemize}
\raggedright
  \item \emph{Say it:} \emph{naitiktā}
  \item \emph{Say it:} \emph{naitiktā}, once more
  \item \emph{Say it:} \emph{vyāpārāthī naitiktā pāḷlī pāhije.} — Ethics should be observed in trade too.
  \item \emph{Recall:} say the Marathi for to fight for, then the Marathi for a value, then say \emph{naitiktā} again
\end{itemize}

\begin{tcolorbox}[breakable,colback=teal!4,colframe=teal!35!black,title={Before you move on}]
What does \mr{नैतिकता} mean? (\textquotedblleft{}Morality; ethics\textquotedblright{}.) Say it once more.
\end{tcolorbox}
\section[\texorpdfstring{\mr{सचोटी} (sacoṭī)}{सचोटी (sacoṭī)}]{\mr{सचोटी} (sacoṭī) — integrity; honesty}

\label{lesson:MR-C366-sacoti}



\begin{quote}
Before the new one: say the Marathi for a value, then the Marathi for morality.
\end{quote}

\subsection*{You'll want to know: \mr{सचोटी}}
\textbf{\mr{सचोटी}} — \emph{sacoṭī} — \textquotedblleft{}integrity; honesty\textquotedblright{}.

\begin{quote}
\mr{तो} \mr{दुकानदार} \mr{सचोटीने} \mr{व्यवहार} \mr{करतो}.
\end{quote}

\begin{quote}
\emph{to dukāndār sacoṭīne vyavhār karto.}
\end{quote}

\begin{quote}
That shopkeeper deals with integrity.
\end{quote}

Feminine.

\mr{सचोटीने}, honestly.

\begin{cousinweb}
\textbf{\mr{सचोटी}} (\emph{sacoṭī}) is a tadbhava: CDIAL lists Marathi \emph{saċoṭī}, \textquotedblleft{}honesty\textquotedblright{}, under Sanskrit \emph{satyavṛtti}, \textquotedblleft{}devotion to truth\textquotedblright{} (CDIAL 13115).
\end{cousinweb}

\begin{grammarlens}[title={What you've built}]
One more word for giving reasons and opinions.
\end{grammarlens}

\subsection*{Your turn}
\begin{itemize}
\raggedright
  \item \emph{Say it:} \emph{sacoṭī}
  \item \emph{Say it:} \emph{sacoṭī}, once more
  \item \emph{Say it:} \emph{to dukāndār sacoṭīne vyavhār karto.} — That shopkeeper deals with integrity.
  \item \emph{Recall:} say the Marathi for a value, then the Marathi for morality, then say \emph{sacoṭī} again
\end{itemize}

\begin{tcolorbox}[breakable,colback=teal!4,colframe=teal!35!black,title={Before you move on}]
What does \mr{सचोटी} mean? (\textquotedblleft{}Integrity; honesty\textquotedblright{}.) Say it once more.
\end{tcolorbox}
\section[\texorpdfstring{\mr{त्याग} (tyāg)}{त्याग (tyāg)}]{\mr{त्याग} (tyāg) — renunciation; a sacrifice}

\label{lesson:MR-C366-tyag}



\begin{quote}
Before the new one: say the Marathi for morality, then the Marathi for integrity.
\end{quote}

\subsection*{You'll want to know: \mr{त्याग}}
\textbf{\mr{त्याग}} — \emph{tyāg} — \textquotedblleft{}renunciation; a sacrifice\textquotedblright{}.

\begin{quote}
\mr{मुलांसाठी} \mr{आईने} \mr{खूप} \mr{त्याग} \mr{केला}.
\end{quote}

\begin{quote}
\emph{mulāṁsāṭhī āīne khūp tyāg kelā.}
\end{quote}

\begin{quote}
The mother made many sacrifices for her children.
\end{quote}

Masculine.

\mr{त्याग} \mr{करणे}, to give up: \mr{सुखाचा} \mr{त्याग}.

\begin{cousinweb}
\textbf{\mr{त्याग}} (\emph{tyāg}) is a tatsama, taken from Sanskrit \emph{tyāga}, \textquotedblleft{}abandonment, liberality\textquotedblright{} (CDIAL 5985).
\end{cousinweb}

\begin{grammarlens}[title={What you've built}]
One more word for giving reasons and opinions.
\end{grammarlens}

\subsection*{Your turn}
\begin{itemize}
\raggedright
  \item \emph{Say it:} \emph{tyāg}
  \item \emph{Say it:} \emph{tyāg}, once more
  \item \emph{Say it:} \emph{mulāṁsāṭhī āīne khūp tyāg kelā.} — The mother made many sacrifices for her children.
  \item \emph{Recall:} say the Marathi for morality, then the Marathi for integrity, then say \emph{tyāg} again
\end{itemize}

\begin{tcolorbox}[breakable,colback=teal!4,colframe=teal!35!black,title={Before you move on}]
What does \mr{त्याग} mean? (\textquotedblleft{}Renunciation; a sacrifice\textquotedblright{}.) Say it once more.
\end{tcolorbox}
\section[\texorpdfstring{\mr{हित} (hit)}{हित (hit)}]{\mr{हित} (hit) — welfare; good}

\label{lesson:MR-C366-hit}



\begin{quote}
Before the new one: say the Marathi for integrity, then the Marathi for renunciation.
\end{quote}

\subsection*{You'll want to know: \mr{हित}}
\textbf{\mr{हित}} — \emph{hit} — \textquotedblleft{}welfare; good\textquotedblright{}.

\begin{quote}
\mr{हा} \mr{निर्णय} \mr{सगळ्यांच्या} \mr{हिताचा} \mr{आहे}.
\end{quote}

\begin{quote}
\emph{hā nirṇay sagaḷyāñcyā hitācā āhe.}
\end{quote}

\begin{quote}
This decision is for everyone's good.
\end{quote}

Neuter.

X \mr{च्या} \mr{हिताचा}, in X's interest.

\begin{cousinweb}
\textbf{\mr{हित}} (\emph{hit}) is a tatsama, taken from Sanskrit \emph{hita}, \textquotedblleft{}beneficial; welfare, good\textquotedblright{} (Monier-Williams).
\end{cousinweb}

\begin{grammarlens}[title={What you've built}]
One more word for giving reasons and opinions.
\end{grammarlens}

\subsection*{Your turn}
\begin{itemize}
\raggedright
  \item \emph{Say it:} \emph{hit}
  \item \emph{Say it:} \emph{hit}, once more
  \item \emph{Say it:} \emph{hā nirṇay sagaḷyāñcyā hitācā āhe.} — This decision is for everyone's good.
  \item \emph{Recall:} say the Marathi for integrity, then the Marathi for renunciation, then say \emph{hit} again
\end{itemize}

\begin{tcolorbox}[breakable,colback=teal!4,colframe=teal!35!black,title={Before you move on}]
What does \mr{हित} mean? (\textquotedblleft{}Welfare; good\textquotedblright{}.) Say it once more.
\end{tcolorbox}
